\section{Displayed formulas and prose}
\label{sec:writing-foundations-05}

Displayed formulas are parts of sentences and receive punctuation.  Each
symbol must be introduced before the display in which it first appears.
Long transformations should use an aligned environment with a justification at
each nontrivial step:
\[
\begin{aligned}
  \neg(P\to Q)
    &\equiv \neg(\neg P\lor Q) && \text{equivalence of }P\to Q
    \text{ with }\neg P\lor Q\\
    &\equiv P\land\neg Q       && \text{De Morgan and double-negation
    equivalences}.
\end{aligned}
\]
The symbol \(\equiv\) in this example is a metalinguistic equivalence sign; it
is not an object-language connective.  The object-language biconditional will
be written \(\leftrightarrow\).

Halmos, Knuth, Larrabee, and Roberts emphasize that rigorous prose must
remain grammatical and that notation should reduce, not replace, explanation
\cite{Halmos1970,KnuthLarrabeeRoberts1989}.  The AMS author guidance provides
the corresponding publication and \LaTeX{} conventions
\cite{AMSAuthorHandbook}.

